\documentclass[12pt]{article}
\usepackage[T1]{fontenc}
\usepackage[utf8]{inputenc}
\usepackage{lmodern}
\usepackage{textcomp}
\usepackage{enumitem}
\usepackage{geometry}
\geometry{margin=1in}
\begin{document}
\begin{center}
Answer Key - Machine Learning - Undergraduate Level Assessment 11 \
\small (Answers are ordered exactly as in the question paper.)
\end{center}

\section*{Multiple Choice}
\begin{enumerate}[label=\arabic*.]
\item \textbf{Answer:} B
\\ \textit{Options:} A) pure; B) not pure; C) useful; D) useless
\\ \textit{Note:} Correct option: B - not pure
\item \textbf{Answer:} B
\\ \textit{Options:} A) Attributes are equally important.; B) Attributes are statistically dependent of one another given the class value.; C) Attributes are statistically independent of one another given the class value.; D) Attributes can be nominal or numeric
\\ \textit{Note:} Correct option: B - Attributes are statistically dependent of one another given the class value.
\item \textbf{Answer:} C
\\ \textit{Options:} A) True, True; B) False, False; C) True, False; D) False, True
\\ \textit{Note:} Correct option: C - True, False
\item \textbf{Answer:} C
\\ \textit{Options:} A) P(H, U, P, W) = P(H) * P(W) * P(P) * P(U); B) P(H, U, P, W) = P(H) * P(W) * P(P | W) * P(W | H, P); C) P(H, U, P, W) = P(H) * P(W) * P(P | W) * P(U | H, P); D) None of the above
\\ \textit{Note:} Correct option: C - P(H, U, P, W) = P(H) * P(W) * P(P | W) * P(U | H, P)
\item \textbf{Answer:} A
\\ \textit{Options:} A) p(y|x, w); B) p(y, x); C) p(w|x, w); D) None of the above
\\ \textit{Note:} Correct option: A - p(y|x, w)
\end{enumerate}

\section*{True / False}
\begin{enumerate}[label=\arabic*.]
\setcounter{enumi}{5}
\item \textbf{Answer:} True
\\ \textit{Options:} A) True; B) False
\\ \textit{Note:} LLM-generated true statement based on MCQ.
\item \textbf{Answer:} True
\\ \textit{Options:} A) True; B) False
\\ \textit{Note:} LLM-generated true statement based on MCQ.
\item \textbf{Answer:} True
\\ \textit{Options:} A) True; B) False
\\ \textit{Note:} LLM-generated true statement based on MCQ.
\item \textbf{Answer:} True
\\ \textit{Options:} A) True; B) False
\\ \textit{Note:} LLM-generated true statement based on MCQ.
\item \textbf{Answer:} False
\\ \textit{Options:} A) True; B) False
\\ \textit{Note:} LLM-generated false statement. Correct answer: 'It discovers causal relationships'.
\end{enumerate}

\section*{Long Form Response}
\begin{enumerate}[label=\arabic*.]
\setcounter{enumi}{10}
\item \textbf{Answer:} def adjac(ele, sub = []): | def get\_coordinates(test\_tup): | return (res)
\\ \textit{Note:} Students should provide a detailed explanation referencing algorithm steps.
\item \textbf{Answer:} def count\_odd(array\_nums): | return count\_odd
\\ \textit{Note:} Students should provide a detailed explanation referencing algorithm steps.
\end{enumerate}

\vfill
\noindent\textit{End of Answer Key}
\end{document}